\documentclass{../../../unipd-notes}
\unipdsetup{course = {Fondamenti di Ingegneria Informatica},short-course = {Fondamenti di Ingegneria},document-type = {Appunti delle lezioni}}
\begin{document}
\chapter{Elenchi}
\section{Puntati e numerati}
\begin{itemize}
  \item unità aritmetico-logica per l'esecuzione delle operazioni;
  \item unità di controllo, composta da:
    \begin{itemize}
      \item registro dell'istruzione corrente;
      \item logica di decodifica, che comprende:
        \begin{itemize}
          \item riconoscimento dell'opcode;
          \item generazione dei segnali di controllo.
        \end{itemize}
    \end{itemize}
\end{itemize}
\begin{enumerate}
  \item prelevare l'istruzione dalla memoria;
  \item decodificare l'opcode e gli operandi:
    \begin{enumerate}
      \item identificare il formato dell'istruzione;
      \item interpretare gli operandi:
        \begin{enumerate}
          \item leggere i registri sorgente;
          \item estendere gli operandi immediati.
        \end{enumerate}
    \end{enumerate}
  \item eseguire l'operazione e aggiornare lo stato.
\end{enumerate}
\section{Descrizioni e passi}
\begin{description}
  \item[Bit] Una cifra binaria.
  \item[Byte] Una sequenza di otto bit.
\end{description}
\begin{steps}
  \item Tradurre l'indirizzo virtuale mediante la tabella delle pagine.
  \item Accedere al frame fisico e verificare i permessi.
    \begin{steps}
      \item Controllare il bit di presenza.
      \item Verificare i permessi di lettura o scrittura.
    \end{steps}
\end{steps}
\end{document}
